%% sections/07_methodology.tex

\section{실험 방법}
\label{sec:methodology}

\subsection{하드웨어 및 소프트웨어 환경}

본 연구는 두 측정 플랫폼을 사용한다.
\emph{프로덕션 플랫폼}(DGX B200)에서 대규모 실 트레이스 평가---8개 모델 $\times$
4개 방법 $\times$ 7개 조건의 처리량 oracle 및 llm-d 라우팅 비교---를 수행하고,
\emph{개발 플랫폼}(H100$\times$8)에서 CPU 슬랙 하베스팅 lever 및 소형 모델 회귀
cross-validation 등 통제된 microstudy 를 수행한다.

\begin{table}[t]
\centering
\caption{실험 환경. 대규모 실 트레이스 평가(\S\ref{subsec:res-b200})는 DGX B200,
  CPU lever 및 회귀 microstudy 는 개발 플랫폼(H100$\times$8)에서 수행.}
\label{tbl:setup}
\footnotesize
\begin{tabular}{@{}ll@{}}
\toprule
\textbf{항목} & \textbf{사양} \\
\midrule
\multicolumn{2}{@{}l}{\textit{프로덕션 플랫폼 (대규모 실 트레이스 평가)}} \\
GPU & NVIDIA B200 $\times$ 8 (183\,GB HBM3e, sm\_100) \\
CPU & Intel Xeon Platinum 8570, 듀얼 NUMA \\
    & 224 스레드 (2 소켓), DDR5 \\
메모리 & 듀얼 NUMA, $\sim$2\,TB DRAM \\
\midrule
\multicolumn{2}{@{}l}{\textit{개발 플랫폼 (CPU lever $\cdot$ 회귀 microstudy)}} \\
GPU & NVIDIA H100 SXM5 $\times$ 8 (80\,GB HBM3) \\
CPU & Intel Xeon SPR-class, 듀얼 NUMA \\
    & 노드당 56 물리 코어, NUMA distance 2.1$\times$ \\
\midrule
\multicolumn{2}{@{}l}{\textit{공통 소프트웨어}} \\
vLLM & 본 연구 포크 (B200 는 sm\_100 재빌드) \\
SuffixDecoding & arctic-inference (lazy import) \\
라우팅 비교 & llm-d Gateway (Istio $+$ InferencePool $+$ EPP) \\
분류기 & FastAPI $+$ ProcessPool 16 worker \\
측정 도구 & vLLM benchmark\_serving, UtilSampler \\
          & (nvidia-smi $+$ /proc/stat 동기 샘플링) \\
\midrule
모델 & \multicolumn{1}{p{0.62\columnwidth}@{}}{\raggedright
  Qwen-2.5-7B/32B/72B, Llama-3.1-8B/70B,
  DeepSeek-R1-Distill-Qwen-7B/32B, \mbox{-Llama-70B}
  (회귀 microstudy: opt-125m, starcoder2-3b, Qwen-0.5B/1.5B/7B, Phi-3-14B)} \\
\bottomrule
\end{tabular}
\end{table}

\subsection{측정 시나리오}

본 연구는 세 가지 측정 시나리오를 사용한다.

\textbf{시나리오 1 — 대규모 실 트레이스 oracle (B200).}
8개 모델 각각을 단일 vLLM 인스턴스(TP=8, 7B 급은 TP=4)로 기동하고,
4개 방법(vanilla / suffix / ngram / llm-d)으로 동일 prompt 집합을
\emph{조건별}로 측정한다 (concurrency 32, \texttt{max\_tokens}=8{,}192, streaming).
방법별 per-prompt 처리량으로부터 (모델, 조건) 단위의 \emph{oracle}
(=최적 방법)을 산출하여 게이팅 상한과 방법 선택 효과를 정량화한다
(\S\ref{subsec:res-b200}).

\textbf{시나리오 2 — llm-d 라우팅 비교 (B200).}
vanilla 와 suffix 두 백엔드를 \emph{동일 GPU 예산}으로 분할
(8-GPU 모델 = 2$\times$TP4, 4-GPU 모델 = 2$\times$TP2)한 뒤
llm-d Gateway(Istio $+$ InferencePool $+$ EPP)가 load/prefix-aware
라우팅을 수행한다. kind 기반 Kubernetes 클러스터에 vLLM pod 를 배치하고
동일 prompt 집합을 동일 부하(conc=32)로 인가하여 \emph{시스템 전체 처리량}
관점에서 단일 백엔드(vanilla/suffix) 대비 비교한다.

\textbf{시나리오 3 — End-to-end dual-backend (CPU lever, H100).}
standard backend(GPU 0--3)와 spec-decode backend(GPU 4--7)를
동시에 기동하고, \algname\ 분류기가 요청을 concurrency 32로 분배한다.
CPU 슬랙 하베스팅 lever 의 ON/OFF 쌍 측정 및 게이팅 분해에 사용한다.

\subsection{워크로드 정의}

\textbf{실 트레이스 corpus (시나리오 1, 2).}
합성 데이터 대신 6종의 공개 실 서빙/벤치마크 trace 를 prompt 소스로 사용한다:
대화 계열 \textbf{ShareGPT}, \textbf{WildChat}, \textbf{LMSYS-Chat-1M}과
코드 계열 \textbf{SWE-bench}, \textbf{MBPP}, \textbf{HumanEval}.
각 corpus 에서 마지막 user turn(대화) / problem statement(SWE-bench) /
prompt(HumanEval/MBPP)를 추출하고, sha1 정규화 dedup 및 길이 필터
($32 \le \text{tok} \le 4{,}096$)를 적용하여 총 2{,}159 prompt 를 샘플링한다
(대화 3종 각 500, SWE-bench 297, MBPP 198, HumanEval 164).
\textbf{mix} 조건은 6 corpus 를 교차 셔플하여 단일 서버에 이질적 트래픽이
혼재하는 운영 상황을 모사한다.

\textbf{게이팅 라벨 분류.}
워크로드 레벨 게이팅은 corpus 출처가 아니라 요청 내용을
sonnet(창작)/chat(대화)/code(코드)로 분류하여 방법을 선택한다.
즉 corpus 는 prompt \emph{소스}이고, 게이팅 라벨은 분류기가 부여한다
(\S\ref{sec:ceres}). 이로써 합성 분류에 의존하던 이전 평가의
in-distribution 편향을 제거하고 실제 트래픽 분포에서 게이팅을 검증한다.

\subsection{기준선 및 공정 비교}

모든 측정에서 공정 비교(fair comparison)를 위해
baseline과 spec-decode가 동일한 \texttt{gpu\_memory\_utilization=0.80} 및
동일한 벤치마크 래퍼(\texttt{run\_workload\_gen.py})를 사용한다.
이전 측정의 래퍼 편향($+64.7\%$ wrapper noise)을
제거한 fair baseline을 적용한다.

\subsection{CPU 슬랙 하베스팅 측정 방법}

lever 단독 기여 $\Delta(L)$는 \algname-gated ON/OFF 쌍 측정으로 분리한다.
핵심 lever(NUMA pin, softmax aux)에 대해 3-run 반복으로
multi-run mean 및 warm-only(run 2, 3 평균)를 보고한다.
신뢰도 기준: $|\Delta| < 3\%$인 1-run 신호는 multi-run binding 없이 채택하지 않는다.

% TODO: CERES 완전 통합 설정(NUMA pin + branchy 100ms + AGSD-gated)의
%       end-to-end 3-run 측정 필요.
%       현재는 각 lever를 독립 측정 후 Property 1로 결합 추정.
